\chapter{前馈神经网络}
\noindent\key{本章速览}\quad 神经元 = 线性运算 $z=\bm w^\top\bm x+b$ + 非线性激活 $a=f(z)$；逐层堆叠成前馈网络 $\bm z^{(l)}=W^{(l)}\bm a^{(l-1)}+\bm b^{(l)}$，$\bm a^{(l)}=f_l(\bm z^{(l)})$；单隐藏层即可任意逼近连续函数（通用近似定理）；最后一层接 Logistic/Softmax 做分类，用结构化风险 + 梯度下降学参数。
\key{记号}\quad 输入层为第 0 层，$\bm a^{(0)}=\bm x$；$L$ 为层数（不含输入层）；$M_l$ 为第 $l$ 层神经元数；$\sigma(\cdot)$ 为 Logistic 函数；$\Phi,\varphi$ 为标准正态分布函数与密度；$\mathbf 1\{\cdot\}$ 为指示函数。

% ======================================================================
\section{神经网络概述：连接主义、PDP 网络与误差反向传播}

\subsection{连接主义模型；PDP 网络的三个主要特性}
\begin{itemize}
\item 神经网络最早是作为一种主要的\term{连接主义}{connectionism}模型：智能来自大量简单单元及其连接，信息由许多单元的激活模式共同表示，学习 = 调整连接强度（权重）。对立路线：符号主义（显式规则 + 逻辑推理）。\textit{Connectionism models intelligence as networks of simple interconnected units; knowledge is stored in the connection weights.}
\item 80 年代后期最流行的连接主义模型：\term{分布式并行处理}{Parallel Distributed Processing, PDP}网络，\key{三个主要特性}：(1) 信息表示是\key{分布式的（非局部的）}；(2) 记忆和知识存储在\key{单元之间的连接}上；(3) 通过\key{逐渐改变单元之间的连接强度}来学习新知识。对应前馈网络：表示 = 活性值 $\bm a^{(l)}$，知识 = 权重 $W^{(l)},\bm b^{(l)}$，学习 = 梯度下降修改 $W^{(l)}$。
\item 局部表示（one-hot，一个单元一个概念，3 个单元只能表示 3 种颜色）vs 分布式表示（多个单元取值组合，如 RGB）：后者更紧凑、相近概念向量相近、单元损坏不致完全丢失信息（词向量即分布式表示）。
\end{itemize}

\subsection{误差反向传播：神经网络走向主流的关键}
早期感知器只有一层可学习权重、表达能力有限；多层网络缺少训练方法。\term{误差反向传播}{error backpropagation, BP}（Rumelhart、Hinton、Williams 1986）高效计算多层网络梯度；引入 BP 改进学习能力之后，神经网络越来越多地应用于各种机器学习任务。

% ======================================================================
\section{神经元}

\subsection{生物神经元的七个组成部分；兴奋与抑制两种状态}
按信号流向：\key{树突}（Dendrite，接收信号）$\to$ \key{树突分支}（Dendritic branches，扩大接收面）$\to$ \key{细胞体}（Cell body，整合输入）$\to$ \key{轴突丘}（Axon hillock，超阈值时触发冲动）$\to$ \key{轴突}（Axon，长距离传出）$\to$ \key{轴突终末分支}（Telodendria，末端分叉）$\to$ \key{突触终端}（Synaptic terminals，传给下一个细胞）。\key{单个神经细胞只有两种状态：兴奋和抑制}（汇总输入超阈值则兴奋，否则抑制）；激活函数是这种“开/关”行为的光滑化或分段线性化。

\subsection{人工神经元的五个部分及其生物学对应}
\key{输入} $x_1,\dots,x_d$ $\leftrightarrow$ 树突接收的信号；\key{权重} $w_1,\dots,w_d$ $\leftrightarrow$ 突触连接强度；\key{偏置} $b$ $\leftrightarrow$ 触发阈值的调整项；\key{求和单元} $z=\sum_iw_ix_i+b$ $\leftrightarrow$ 细胞体的信号整合；\key{激活函数与输出} $a=f(z)$ $\leftrightarrow$ 兴奋/抑制状态。

\subsection{数学模型：净输入 \texorpdfstring{$z$}{z} 与活性值 \texorpdfstring{$a$}{a}；“线性运算 + 非线性转换”}
$z=\sum_{i=1}^dw_ix_i+b=\bm w^\top\bm x+b$，$a=f(z)$。$z$：\term{净输入}{net input}（净活性值）；$a$：\term{活性值}{activation}（输出）；$f$：\term{激活函数}{activation function}。
\begin{itemize}
\item \key{偏置 = 常数输入 1 的权重}：$\tilde{\bm x}=(1,x_1,\dots,x_d)^\top$，$\tilde{\bm w}=(b,w_1,\dots,w_d)^\top$，$z=\tilde{\bm w}^\top\tilde{\bm x}$。
\item \key{本质}：“\key{线性运算 + 非线性转换}”的信号处理器——$z=\bm w^\top\bm x+b$ 为仿射运算，$a=f(z)$ 为非线性转换。Logistic 回归就是 $f=\sigma$ 的单个神经元。示意：$\bm x=(2,1,0.5)^\top$，$\bm w=(0.5,-1,2)^\top$，$b=-0.5\Rightarrow z=0.5$，$\sigma(z)\approx0.6225$，$\tanh z\approx0.4621$，$\mathrm{ReLU}(z)=0.5$。
\end{itemize}

% ======================================================================
\section{激活函数}

\subsection{激活函数应满足的四条性质}
(1) \key{连续并可导（允许少数点上不可导）的非线性函数}——可导才能直接用数值优化方法学习参数；(2) \key{激活函数及其导函数尽可能简单}——提高计算效率，如 Sigmoid $f'=f(1-f)$；(3) \key{导函数的值域在合适区间内}，不能太大也不能太小——否则影响训练效率和稳定性（太小梯度消失、太大梯度爆炸）；(4) \key{单调递增}——梯度方向更稳定。\textbf{易错} 这是经典准则而非硬性规定：Swish、GELU 连续可导、导数有界，但\key{非单调}。

\subsection{为什么必须非线性：多层线性网络退化为单层仿射模型}
若各层 $f(z)=z$：$\bm a^{(2)}=W^{(2)}W^{(1)}\bm x+W^{(2)}\bm b^{(1)}+\bm b^{(2)}$，归纳得
\[
\bm a^{(L)}=\underbrace{W^{(L)}\cdots W^{(1)}}_{\tilde W}\bm x+\underbrace{\textstyle\sum_{l=1}^{L}W^{(L)}\cdots W^{(l+1)}\bm b^{(l)}}_{\tilde{\bm b}} \why{仿射的复合仍是仿射；$l=L$ 项连乘为 $I$}
\]
$\Rightarrow$ 叠多少层都等价于 $\tilde W\bm x+\tilde{\bm b}$，只能给出线性决策边界（学不会异或）。

\subsection{Logistic（Sigmoid）函数：定义、值域、\texorpdfstring{$\sigma'=\sigma(1-\sigma)\le\frac14$}{sigma'=sigma(1-sigma)<=1/4}}
$\sigma(x)=\frac{1}{1+\exp(-x)}$，$\R\to(0,1)$，$\sigma(0)=0.5$，$\sigma(-x)=1-\sigma(x)$；Logistic 函数 = Sigmoid 函数。
\begin{align*}
\sigma'(x)&=\frac{e^{-x}}{(1+e^{-x})^2}=\sigma(x)\cdot\frac{(1+e^{-x})-1}{1+e^{-x}}=\sigma(x)\big(1-\sigma(x)\big)\le\Big(\frac{\sigma+1-\sigma}{2}\Big)^2=\frac14 \why{链式法则；加 1 减 1；均值不等式}
\end{align*}
等号仅在 $x=0$；$0<\sigma'\le\frac14$。数值：$\sigma(2)\approx0.8808$，$\sigma'(2)\approx0.105$，$\sigma'(5)\approx0.0066$。

\subsection{Tanh 函数：定义、\texorpdfstring{$\tanh'=1-\tanh^2\le1$}{tanh'=1-tanh^2<=1}、\texorpdfstring{$\tanh(x)=2\sigma(2x)-1$}{tanh(x)=2sigma(2x)-1}}
$\tanh(x)=\frac{e^x-e^{-x}}{e^x+e^{-x}}$，$\R\to(-1,1)$，$\tanh(0)=0$。记 $u=e^x-e^{-x}$，$v=e^x+e^{-x}$（$u'=v,v'=u$）：$\tanh'=\frac{v^2-u^2}{v^2}=1-\tanh^2(x)\in(0,1]$，最大值 1（$x=0$），$\tanh'(2)\approx0.071$。
关系：$2\sigma(2x)-1=\frac{1-e^{-2x}}{1+e^{-2x}}=\frac{e^x-e^{-x}}{e^x+e^{-x}}=\tanh(x)$ \whyin{通分；同乘 $e^x$}，即 Logistic 横向压缩一半、纵向拉伸 2 倍、下移 1；$\tanh'(x)=4\sigma'(2x)\le1$；$\tanh(1)=2\sigma(2)-1=0.7616$。

\subsection{“饱和”的含义；零中心化与偏置偏移（bias shift）}
\begin{itemize}
\item \term{饱和}{saturating}：$\lim_{x\to-\infty}f'(x)=0$ 为左饱和，$\lim_{x\to+\infty}f'(x)=0$ 为右饱和，两者皆成立为（两端）饱和。Logistic、Tanh 是\key{两端饱和}的经典饱和型激活：$|x|$ 稍大时导数几乎为 0 $\Rightarrow$ 梯度消失的根源。
\item \term{零中心化}{zero-centered}：输出关于 0 对称、有正有负。\key{Tanh 零中心化；Logistic 输出恒大于 0}。\textbf{易错}“Sigmoid 是零中心化的”——错。
\item \term{偏置偏移}{bias shift}：非零中心化的输出使后一层神经元的输入发生偏置偏移，使梯度下降收敛变慢。原因：后一层 $z=\bm w^\top\bm a+b$，$a_j>0$ 时 $\frac{\partial L}{\partial w_j}=\frac{\partial L}{\partial z}a_j$ 的符号对所有 $j$ 相同，一次更新各分量同号，最优方向“有增有减”时只能锯齿形逼近；且 $\E[z]=\bm w^\top\E[\bm a]+b$ 整体偏离 0。
\end{itemize}

\subsection{ReLU 系列：ReLU、Leaky ReLU、PReLU、ELU、Softplus}
（公式与导数见 §\ref{sec:D03a-acttab} 总表。）
\begin{itemize}
\item \term{ReLU}{Rectified Linear Unit}：$\max(0,x)$，导数 $\mathbf 1\{x>0\}$；$x=0$ 处左导 0、右导 1，\key{不可导}，约定 $f'(0)=0$（或取 $[0,1]$ 中任一值，次梯度）。正半轴导数恒为 1（非饱和），负半轴导数 0（严格说左饱和）。
\item \term{Leaky ReLU}{带泄露的 ReLU}：负半轴 $\gamma x$，$\gamma$ \key{固定}小常数（如 0.01）；负半轴导数 $\gamma>0$，误差信号不被截断。
\item \term{PReLU}{Parametric ReLU}：第 $i$ 个神经元斜率 $\gamma_i$ \key{可学习}，$\partial f/\partial\gamma_i=\min(0,x)$，与权重一起梯度下降；与 Leaky ReLU 唯一区别是 $\gamma$ 固定还是学出来。
\item \term{ELU}{Exponential Linear Unit}：负半轴 $\gamma(e^x-1)$，$\gamma>0$；负半轴 $f'=\gamma e^x=f+\gamma$；$x\to-\infty$ 时\key{平滑趋于 $-\gamma$}；输出均值更接近 0，缓解偏置偏移；仅 $\gamma=1$ 时 0 点可导。
\item \term{Softplus}{Softplus}：$\log(1+e^x)$，$\mathrm{softplus}(0)=\ln2\approx0.693$，导数 $=\sigma(x)$。ReLU 的\key{光滑版本}（$x\to+\infty$ 时 $\to x$，$x\to-\infty$ 时 $\to0$），但负半轴导数趋 0（仍饱和）、无稀疏性。
\item 示意（$\gamma=0.1$，$x=-2$，值/导）：ReLU $0/0$；Leaky ReLU $-0.2/0.1$；ELU $-0.0865/0.0135$。
\end{itemize}

\subsection{ReLU 系列的三个核心优势；死亡 ReLU 问题}
\key{优势}：(1) \key{计算更高效}（比较取最大，无指数，导数 0 或 1）；(2) \key{生物学合理性}（单侧抑制、宽兴奋边界）；(3) \key{一定程度上缓解梯度消失}（正半轴导数恒为 1，不像 Sigmoid 每层至少缩小到 $\frac14$）。\tagext\ 部分输出恰为 0，激活稀疏。
\term{死亡 ReLU 问题}{dying ReLU problem}：输入长期处于负区间时，ReLU 输出恒为 0、导数也为 0，反向传播时参数无法更新，“永久抑制”（死亡）；常见诱因为学习率过大把偏置推成很大的负数。\key{解决}：用 Leaky ReLU、PReLU、ELU 等变体，让负区间有非零输出。

\subsection{Swish：定义 \texorpdfstring{$x\sigma(\beta x)$}{x sigma(beta x)}、参数 \texorpdfstring{$\beta$}{beta} 的作用、导数、四个特点}
$\mathrm{swish}(x)=x\,\sigma(\beta x)$，$\beta$ 可调节（也可学习），控制 Sigmoid 的\key{陡峭程度}；$\sigma(\beta x)$ 为“软门控”。$\beta=0$：$x/2$（直线）；$\beta\to\infty$：趋于 ReLU；中间 $\beta$：负半轴有谷底、\key{非单调}（$\beta=1$ 时最小约 $-0.278$，在 $x\approx-1.278$）。
\begin{itemize}
\item 导数：$\mathrm{swish}'(x)=\sigma(\beta x)+\beta x\,\sigma(\beta x)\big(1-\sigma(\beta x)\big)=\beta\,\mathrm{swish}(x)+\sigma(\beta x)\big(1-\beta\,\mathrm{swish}(x)\big)$；$\beta=1$ 时取值约 $[-0.100,1.100]$。示意 $\mathrm{swish}(-1)=-0.2689$，$\mathrm{swish}'(-1)\approx0.0723$。
\item \key{四个特点}：(1) \key{平滑可导}（全局连续可导，梯度更稳定）；(2) \key{非单调}（与 ReLU、Sigmoid、Tanh 的最大区别，更灵活拟合复杂分布）；(3) \key{缓解死亡 ReLU}（负区间输出小的负数，参数仍可更新）；(4) \key{兼具 ReLU 与 Sigmoid 的优势}（$\beta$ 较大接近 ReLU 的高效，$\beta$ 适中保留 Sigmoid 的平滑）。
\end{itemize}

\subsection{GELU：定义 \texorpdfstring{$x\Phi(x)$}{x Phi(x)}、导数、“概率权重”与两个近似式}
\term{GELU}{Gaussian Error Linear Unit，高斯误差线性单元}：在 Swish 基础上优化，用高斯分布的累积分布函数替代 Sigmoid：$\mathrm{GELU}(x)=x\,P(X\le x)$，$X\sim N(\mu,\sigma^2)$，$\mu,\sigma$ 为超参数，一般 $\mu=0,\sigma=1$，即 $\mathrm{GELU}(x)=x\Phi(x)$，$\Phi(x)=\frac12[1+\operatorname{erf}(x/\sqrt2)]$（erf：高斯误差函数）。\textbf{易错} $N(\mu,\sigma^2)$ 中 $\sigma$ 是标准差，近似式 $x\sigma(1.702x)$ 中 $\sigma$ 是 Logistic 函数。
\begin{itemize}
\item 导数 $\mathrm{GELU}'(x)=\Phi(x)+x\varphi(x)$，取值约 $[-0.129,1.129]$（极值于 $x=\pm\sqrt2$）；处处光滑、负半轴非单调，最小约 $-0.170$（$x\approx-0.752$）。示意 $\mathrm{GELU}(2)=1.9545$，$\mathrm{GELU}'(2)\approx1.085$。
\item \key{概率权重}：$\Phi(x)$ 为标准正态变量不超过 $x$ 的概率；$x\gg0$ 时输出 $\approx x$，$x\ll0$ 时 $\approx0$；输入越偏离 0，权重越接近 1 或 0。\tagext\ 以概率 $\Phi(x)$ 保留 $x$、否则置 0 的随机门，期望即 $x\Phi(x)$。
\item \key{两个近似}：$\mathrm{GELU}(x)\approx0.5x\big(1+\tanh\big(\sqrt{2/\pi}(x+0.044715x^3)\big)\big)$，$\mathrm{GELU}(x)\approx x\,\sigma(1.702x)$。常数来源：设 $\Phi\approx\frac12(1+\tanh(k(x+cx^3)))$，0 点泰勒匹配 $x$ 系数得 $k=\sqrt{2/\pi}\approx0.798$，匹配 $x^3$ 系数得 $c=\frac{4-\pi}{6\pi}\approx0.0455$（0.044715 为更大区间上数值拟合值）；设 $\Phi\approx\sigma(kx)$，1.702 使 $\max_x|\sigma(kx)-\Phi(x)|$ 最小（误差约 0.0095；仅匹配原点斜率得 1.596）。近似 2 即 $\beta=1.702$ 的 Swish。示意 $\mathrm{GELU}(1)\approx0.8413$（两近似 0.8412、0.8458）。
\end{itemize}

\subsection{GELU 与 Swish 的三个核心区别}
(1) \key{非线性来源不同}：Swish 靠 Sigmoid 的 S 型曲线，GELU 靠高斯分布的累积特性，更贴合图像、文本等自然数据的统计分布；(2) \key{曲线平滑度更优}：GELU 全局光滑，正负区间过渡更平缓，梯度波动更小，更适合超深层网络（几十层的 Transformer）；(3) \key{自适应性更强}：输出受“概率权重”调控，保留 ReLU 的高效性，且\key{不用手动或学习调节 $\beta$}。

\subsection{如何选择激活函数（隐藏层与输出层）}
\key{浅层（$\le3$ 层）}：Sigmoid（二分类输出层）、Tanh，计算简单、不用复杂参数调控；\key{普通深层（CNN、简单 DNN）}：优先 ReLU（高效）、Leaky ReLU（缓解死亡神经元）；\key{大型超深层（Transformer、预训练模型）}：首选 GELU（梯度稳、拟合好），次选 Swish（需算力支撑）。\key{核心原则}：平衡“计算效率、梯度稳定性、数据适配性”，优先选没有额外参数、梯度特性好的。\key{输出层由任务定}：回归用恒等，二分类用 Logistic，多分类用 Softmax。

\subsection{常见激活函数及其导数总表（含比较要点）}\label{sec:D03a-acttab}
\begin{center}\small\renewcommand{\arraystretch}{1.2}
\begin{tabularx}{\linewidth}{@{}p{1.55cm}p{4.0cm}p{4.3cm}p{1.55cm}L@{}}
\toprule
名称 & $f(x)$ & $f'(x)$ & 值域 & 饱和/零中心/主要问题 \\
\midrule
Logistic & $\frac1{1+\exp(-x)}$ & $f(1-f)\in(0,\frac14]$ & $(0,1)$ & 两端饱和；非零中心；梯度消失、偏置偏移 \\
Tanh & $\frac{\exp(x)-\exp(-x)}{\exp(x)+\exp(-x)}$ & $1-f^2\in(0,1]$ & $(-1,1)$ & 两端饱和；零中心；梯度消失较轻 \\
ReLU & $\max(0,x)$ & $\mathbf 1\{x>0\}$ & $[0,\infty)$ & 正半轴不饱和；死亡 ReLU、0 点不可导 \\
Leaky ReLU & $\max(0,x)+\gamma\min(0,x)$ & $\mathbf 1\{x>0\}+\gamma\mathbf 1\{x\le0\}$ & $\R$ & 不饱和；需人为选 $\gamma$ \\
PReLU & $\max(0,x)+\gamma_i\min(0,x)$ & 对 $x$ 同上；对 $\gamma_i$：$\min(0,x)$ & $\R$ & $\gamma_i$ 可学习 \\
ELU & $\max(0,x)+\min(0,\gamma(e^x-1))$ & $\mathbf 1\{x>0\}+\mathbf 1\{x\le0\}\gamma e^x$ & $(-\gamma,\infty)$ & 均值近 0；负半轴趋 $-\gamma$ \\
Softplus & $\log(1+\exp(x))$ & $\frac1{1+\exp(-x)}=\sigma(x)$ & $(0,\infty)$ & ReLU 光滑版；左饱和 \\
Swish & $x\sigma(\beta x)$ & $\sigma(\beta x)+\beta x\sigma(\beta x)(1-\sigma(\beta x))$ & 有下界 & 光滑、非单调；$\beta$ 需调/学 \\
GELU & $x\Phi(x)$ & $\Phi(x)+x\varphi(x)$ & $\approx[-0.17,\infty)$ & 光滑、非单调；无额外参数 \\
\bottomrule
\end{tabularx}
\end{center}
\textbf{易错} 每层 Logistic 导数至多 $\frac14$，5 层连乘至多 $(\frac14)^5\approx0.001$（梯度消失）；Swish、GELU 导数可略大于 1。

% ======================================================================
\section{人工神经网络：三要素与三种网络结构}

\subsection{三要素：激活规则、拓扑结构、学习算法}
\term{人工神经网络}{artificial neural network, ANN}由大量神经元及其\key{有向连接}构成，从三方面刻画：(1) \key{神经元的激活规则}——输入到输出的映射，为非线性函数；(2) \key{网络的拓扑结构}——神经元间的连接关系，如深层前馈网络、CNN、RNN；(3) \key{学习算法}——用训练数据学参数，如反向传播、优化器（梯度下降、Adam）。三要素任一改变即为不同模型。

\subsection{三种网络结构：前馈网络、记忆网络、图网络}
由许多神经元组成的信息处理网络具有\key{并行分布结构}；示意图中\key{圆形节点表示一个神经元，方形节点表示一组神经元}。
\begin{center}\small\renewcommand{\arraystretch}{1.1}
\begin{tabularx}{\linewidth}{@{}p{3.1cm}Lp{4.6cm}@{}}
\toprule
结构 & 连接方式 & 典型模型（适合数据） \\
\midrule
前馈网络 & 分层，信息只从前一层流向后一层，无反馈，有向无环图 & 全连接网络、CNN（独立样本特征表、图像） \\
记忆网络（反馈网络） & 还接收自己的历史信息（自环），有记忆，状态随时间变化 & RNN、Hopfield 网络（时间序列、文本） \\
图网络 & 定义在图结构数据上，节点（一个或一组神经元）连接有向或无向，接收相邻节点信息 & 图神经网络 GNN（关系网络） \\
\bottomrule
\end{tabularx}
\end{center}

% ======================================================================
\section{前馈神经网络的结构}

\subsection{定义与三个特征（全连接神经网络、多层感知器 MLP）}
\term{前馈神经网络}{feedforward neural network, FNN}，也称\key{全连接神经网络}、\key{多层感知器}（MLP）：(1) 各神经元分属不同的层，\key{层内无连接}；(2) \key{相邻两层}神经元\key{全部两两连接}；(3) 整个网络\key{无反馈}，信号从输入层向输出层单向传播，可用\key{有向无环图}表示。\textbf{易错} MLP 用连续可导激活，不是阶跃函数感知器的简单叠加；三名同义。

\subsection{输入层、隐藏层、输出层；记号与维度}
\key{输入层}（第 0 层，接收 $\bm x$，不计算）$\to$ \key{隐藏层}（输出不直接观测，逐层变换特征）$\to$ \key{输出层}（给出 $\hat{\bm y}$）。记号：$L$ 层数；$M_l$ 第 $l$ 层神经元个数（也记 $m^{(l)}$）；$f_l(\cdot)$ 第 $l$ 层激活函数（逐元素）；$W^{(l)}\in\R^{M_l\times M_{l-1}}$、$\bm b^{(l)}\in\R^{M_l}$ 为第 $l-1$ 层到第 $l$ 层的权重矩阵与偏置；$\bm z^{(l)}\in\R^{M_l}$ 净输入（净活性值），$\bm a^{(l)}\in\R^{M_l}$ 输出（活性值）。
\textbf{易错} ① 输入层\key{不计入} $L$：4--5--3--1 网络 $L=3$，$M_0=4,\dots,M_3=1$（口语“四层网络”，作答写明约定）；② $W^{(l)}$ 行数 = 本层 $M_l$、列数 = 上层 $M_{l-1}$，$w^{(l)}_{ij}$ 为“上层第 $j$ 个 $\to$ 本层第 $i$ 个”，第 $i$ 行即本层第 $i$ 个神经元的 $\bm w^\top$。

\subsection{逐层前向传播公式}
\[
\bm z^{(l)}=W^{(l)}\bm a^{(l-1)}+\bm b^{(l)},\qquad \bm a^{(l)}=f_l\big(\bm z^{(l)}\big),\qquad l=1,\dots,L .
\]
\key{前向传播}（forward propagation）：$\bm x=\bm a^{(0)}\to\bm z^{(1)}\to\bm a^{(1)}\to\cdots\to\bm a^{(L-1)}\to\bm z^{(L)}\to\bm a^{(L)}=\phi(\bm x;W,\bm b)$（整个网络是以全部 $W,\bm b$ 为参数的复合函数）。分量 $z_i^{(l)}=\sum_{j}w^{(l)}_{ij}a^{(l-1)}_j+b^{(l)}_i$；维度 $(M_l\times M_{l-1})(M_{l-1}\times1)=M_l\times1$。示意（3--2--1，隐藏 ReLU、输出恒等）：$\bm x=(1,2,-1)^\top$，$W^{(1)}=\begin{psmallmatrix}1&0.5&-1\\-0.5&1&2\end{psmallmatrix}$，$\bm b^{(1)}=(-1,-1)^\top\Rightarrow\bm z^{(1)}=(2,-1.5)^\top$，$\bm a^{(1)}=(2,0)^\top$；$W^{(2)}=(1,-2)$，$b^{(2)}=0.5\Rightarrow\hat y=2.5$。

\subsection{参数个数 \texorpdfstring{$P=\sum_l M_l(M_{l-1}+1)$}{P=sum M_l(M_{l-1}+1)}；深层网络的层次化特征}
第 $l$ 层权重 $M_lM_{l-1}$ + 偏置 $M_l$。4--5--3--1：$25+18+4=47$；784--100--50--10：$78500+5050+510=84060$。
\key{深层前馈网络}：层越深特征越抽象——像素 $\to$ \key{边缘}（隐藏层 1）$\to$ \key{边缘的组合/部件}（眼、鼻，隐藏层 2）$\to$ \key{物体模型}（整张人脸，隐藏层 3）；层层组合出\term{层次化特征}{hierarchical features}，即“深度”的价值。

% ======================================================================
\section{通用近似定理}

\subsection{定理陈述（Cybenko 1989；Hornik 等 1989）与挤压函数}
\term{挤压函数}{squashing function}：非常数、有界、单调递增的连续函数（Logistic、Tanh）。
\key{定理}：令 $\phi(\cdot)$ 为非常数、有界、单调递增的连续函数，$\mathcal I_d=[0,1]^d$。对任意 $f\in C(\mathcal I_d)$ 和 $\epsilon>0$，存在整数 $m$、$v_i,b_i\in\R$、$\bm w_i\in\R^d$，使
\[
F(\bm x)=\sum_{i=1}^m v_i\,\phi\big(\bm w_i^\top\bm x+b_i\big)\quad\text{满足}\quad |F(\bm x)-f(\bm x)|<\epsilon,\ \ \forall\bm x\in\mathcal I_d .
\]
\textit{A one-hidden-layer network with a squashing activation and enough hidden units can uniformly approximate any continuous function on a compact set to any accuracy.} Cybenko 对 Sigmoid 型函数证明；Hornik, Stinchcombe \& White 证明多层前馈网络是通用逼近器。

\subsection{网络解读：一个挤压型隐藏层 + 线性输出层}
具有\key{线性输出层}和\key{至少一个使用“挤压”性质激活函数的隐藏层}的前馈网络，只要\key{隐藏层神经元数量足够}，可以任意精度近似 $\R^d$ 中\key{有界闭集}上的任何连续函数。对照：$W^{(1)}$ 第 $i$ 行为 $\bm w_i^\top$，$b^{(1)}_i=b_i$；输出 1 个神经元、恒等激活，$W^{(2)}=(v_1,\dots,v_m)$。结论是\key{一致逼近}。

\subsection{直观：用陡峭的 Sigmoid 拼出“台阶”；一维构造性证明骨架}
$\sigma(w(x-t))$ 在 $w$ 很大时近似 $x=t$ 处的阶跃：位置 $t=-b/w$，陡峭度 $w$，高度由输出权重 $v$ 定；台阶越多拼得越细，函数变化越剧烈需要的隐藏神经元越多（如 1--8--1 网络逼近 $0.4\sin2\pi x+0.5x$，最大误差约 0.098）。
证明（$d=1$，$\phi=\sigma$，$H(u)=\mathbf 1\{u\ge0\}$）：
\begin{enumerate}
\item \key{一致连续}：存在 $K$，$|x-x'|\le\frac1K\Rightarrow|f(x)-f(x')|<\frac\epsilon4$。网格 $x_k=\frac kK$，跳跃位置 $t_k=\frac{k-1/2}{K}$，高度 $c_k=f(x_k)-f(x_{k-1})$，$|c_k|<\frac\epsilon4$。
\item \key{阶梯函数}：$S(x)=f(x_0)+\sum_{k=1}^Kc_kH(x-t_k)=f(x_j)$（望远镜求和，$x_j$ 为最近网格点）$\Rightarrow|S-f|<\frac\epsilon4$。
\item \key{Sigmoid 代替阶跃}：$G(x)=f(x_0)+\sum_kc_k\sigma(\lambda(x-t_k))$；$|\sigma(\lambda u)-H(u)|\le\frac12$，$|u|\ge\frac1{2K}$ 时 $\le\sigma(-\frac\lambda{2K})$，至多一个 $t_k$ 离 $x$ 近 $\Rightarrow|G-S|\le\frac\epsilon8+K\frac\epsilon4\sigma(-\frac\lambda{2K})\le\frac\epsilon4$（$\lambda$ 足够大）；故 $|G-f|<\frac\epsilon2<\epsilon$。
\item \key{化为定理形式}：常数 $f(x_0)=2f(x_0)\sigma(0\cdot x+0)$，$m=K+1$，$v_i=c_i,w_i=\lambda,b_i=-\lambda t_i$。一般挤压函数用 $\frac{\phi(\lambda u)-A}{B-A}$（$A=\phi(-\infty),B=\phi(+\infty)$）代替；多维原证明用泛函分析（Hahn--Banach、Riesz 表示定理）。
\end{enumerate}

\subsection{定理的意义与局限；为什么还要深层网络}
\key{说明了}：表达能力足够强，神经网络可作“万能”函数。\key{没说明}：(1) 只保证参数\key{存在}，不告诉 $m$ 多大（$d$ 维网格构造时随 $d$ 指数增长）；(2) 不保证梯度下降等\key{学习算法能找到}（损失非凸）；(3) 不涉及有限样本的\key{泛化能力}。\textbf{易错}“定理保证梯度下降一定能训练出逼近任意函数的网络”——错。\key{为什么用深层}：逐层组合出层次化特征，往往用少得多的参数达到同样表达能力。\tagext\ Leshno 等（1993）：激活函数只要\key{不是多项式}，单隐藏层就有通用近似性质（ReLU 适用）。

% ======================================================================
\section{神经网络应用到机器学习}

\subsection{特征转换视角 \texorpdfstring{$\hat y=g(\phi(\bm x),\theta)$}{y=g(phi(x),theta)}；二分类：Logistic 回归作最后一层}
神经网络可作“万能”函数，用于复杂的\key{特征转换}或逼近复杂的\key{条件分布}：$\hat y=g\big(\phi(\bm x),\theta\big)$，$\phi$ 为神经网络，$g$ 为分类器，$\bm x$ 为原始输入（像素、词向量），$\theta$ 为整个框架参数（网络权重 + 分类器参数）。$\phi$ 由数据学出（\term{表示学习}{representation learning}），前 $L-1$ 层负责 $\phi$、最后一层负责 $g$，\key{端到端}训练。
若 $g$ 为 Logistic 回归，它可看成网络的最后一层：$\phi(\bm x)=\bm a^{(L-1)}$，$\hat y=\sigma(\bm w^\top\phi(\bm x)+b)$；配交叉熵时 $\partial L/\partial\bm w=(\hat y-y)\phi(\bm x)$，与 Logistic 回归梯度相同，只是特征是学出来的。

\subsection{多分类：Softmax 输出层与交叉熵损失}
最后一层设 $C$ 个神经元，经 Softmax 归一化后作为各类条件概率：
\[
\hat y_c=\frac{\exp(z^{(L)}_c)}{\sum_{k=1}^C\exp(z^{(L)}_k)},\qquad L(\bm y,\hat{\bm y})=-\bm y^\top\log\hat{\bm y}=-\sum_{c=1}^Cy_c\log\hat y_c\quad(\bm y\text{ 为 one-hot，真实类 }c\text{ 时 }L=-\log\hat y_c).
\]
示意 $\bm z^{(L)}=(2,1,0)^\top$，真实第 1 类：$\hat{\bm y}\approx(0.665,0.245,0.090)^\top$，$L\approx0.408$。$C=2$ 时 $\hat y_1=\sigma(z_1-z_2)$，等价于一个 Logistic 神经元。

% ======================================================================
\section{参数学习：结构化风险与梯度下降}

\subsection{结构化风险函数与 Frobenius 范数正则}
\[
\mathcal R(W,\bm b)=\frac1N\sum_{n=1}^N L\big(\bm y^{(n)},\hat{\bm y}^{(n)}\big)+\frac12\lambda\|W\|_F^2,\qquad \|W\|_F^2=\sum_{l=1}^L\sum_{i=1}^{M_l}\sum_{j=1}^{M_{l-1}}\big(w^{(l)}_{ij}\big)^2 .
\]
经验风险 + 正则项，$\lambda>0$ 为超参数；$\|W\|_F^2$ 为全部权重平方和；\key{偏置一般不正则}。正则梯度 $\frac{\partial}{\partial W^{(l)}}\frac12\lambda\|W\|_F^2=\lambda W^{(l)}$（$\frac12$ 与 2 抵消），对 $\bm b^{(l)}$ 为 0。

\subsection{梯度下降的更新规则、学习率；导数的正确定义}
$W^{(l)}\leftarrow W^{(l)}-\alpha\frac{\partial\mathcal R}{\partial W^{(l)}}$，$\bm b^{(l)}\leftarrow\bm b^{(l)}-\alpha\frac{\partial\mathcal R}{\partial\bm b^{(l)}}$；\term{学习率}{learning rate} $\alpha$ 控制步长。算法：(1) 初始化 $\bm w$；(2) 重复：算梯度 $\partial L/\partial\bm w$，更新 $\bm w\leftarrow\bm w-\alpha\,\partial L/\partial\bm w$（全样本为批量 GD，一个/一小批为（小批量）SGD）。示意：经验风险梯度 $(0.2,-0.4)$，$W^{(2)}=(1,2)$，$\lambda=0.1$，$\alpha=0.5\Rightarrow\partial\mathcal R/\partial W^{(2)}=(0.3,-0.2)$，更新为 $(0.85,2.1)$。
导数定义 $\frac{\partial f(w)}{\partial w}=\lim_{\Delta w\to0}\frac{f(w+\Delta w)-f(w)}{\Delta w}$（\textbf{易错} 分子须减 $f(w)$）；多元逐分量求偏导，梯度与参数\key{同形}。

\subsection{为什么沿负梯度方向损失下降；学习率过大会发散}
\begin{itemize}
\item $L(\bm w-\alpha\bm g)=L(\bm w)-\alpha\|\bm g\|^2+o(\alpha)<L(\bm w)$（一阶泰勒，$\bm g\ne\bm 0$，$\alpha$ 充分小）；步长固定时 $\bm g^\top\Delta\bm w\ge-\|\bm g\|\|\Delta\bm w\|$（柯西--施瓦茨），等号当且仅当 $\Delta\bm w\parallel-\bm g$ $\Rightarrow$ 负梯度是局部下降最快方向。
\item $L=\frac c2w^2$：$w_{t+1}=(1-\alpha c)w_t$，$w_t=(1-\alpha c)^tw_0$；\key{收敛当且仅当 $0<\alpha<2/c$}。$c=2$：$\alpha=0.1$ 平稳收敛（$0.8^t$），$\alpha=0.9$ 震荡收敛（$(-0.8)^t$），$\alpha=1.2$ 震荡发散（$(-1.4)^t$）；$\alpha$ 太小收敛很慢。
\end{itemize}

\subsection{梯度下降可能停在局部极小；计算梯度的三种途径}
梯度为 0 的点可能是局部极小、鞍点或平台，梯度下降停在此处，越不过“尖峰”到达更低的全局极小（如 $L=2(\frac{w^4}4-\frac{6.7w^3}3+6.7w^2-7.7w)+7$，$L'=2(w-1)(w-2.2)(w-3.5)$，从 $w_0=-0.1$ 停在局部极小 $w=1$，到不了全局极小 $w=3.5$）。
\key{三种途径}：(1) \key{链式法则} $y=f^5(f^4(f^3(f^2(f^1(x)))))\Rightarrow\frac{\partial y}{\partial x}=\frac{\partial f^1}{\partial x}\frac{\partial f^2}{\partial f^1}\cdots\frac{\partial f^5}{\partial f^4}$（上标为函数编号）——通用，但直接算神经网络梯度很慢；(2) \key{反向传播算法}——根据前馈网络特点设计的高效方法；(3) \key{自动微分}（AD）——PyTorch、TensorFlow 等框架的做法。
